% The Lean build generates the node files.
% Run `lake build :blueprint` before `make pdf` or `make web`.
\input{../../../.lake/build/blueprint/library/Construction}
\input{../../../.lake/build/blueprint/library/Proof}

\chapter{Introduction}

This blueprint describes the three verified security properties of the submitted ArgoMAC garbling scheme for BN254 scalar multiplication from the BABE paper~\cite{babe}.
The scheme garbles the function $f[r](\pi)$ of \cite[Eq.~27]{babe}: it returns $r\pi$ for a point $\pi \in \mathbb{G}_1$ and $0$ otherwise, and it hides the secret scalar $r$.
The Lean project extracts each node from the source with LeanArchitect.
Each node shows the Lean declaration, its statement in the notation of the paper, and its dependencies.
The dependency graph on the web version shows the proof tree of each property.
The roots match the \texttt{perfectCorrectness}, \texttt{adaptivePrivacy}, and \texttt{lamportCompatible} fields of \texttt{Submission.solution}.

\section{Paper correspondence}

Each node names its counterpart in the BaBe paper~\cite{babe}.
The references use BABE v2.0, the version pinned by the Lean sources
(\href{https://github.com/babylonlabs-io/BaBe.latex/tree/e2dcf4d540b2708e13cd21090df759051119a116}{commit \texttt{e2dcf4d}}).
Each chapter restates the items of the BaBe paper before the Lean nodes, with the kind of the item in the paper: theorem, lemma, claim, definition, construction, or algorithm.
Each restatement adapts the wording of the paper to the notation of this blueprint and names the item, for example ``Lemma~14 of the BaBe paper''.
\Cref{tab:paper} lists the paper items and the Lean nodes that formalize them.

\begin{table}[htbp]
\centering
\small
\begin{tabular}{p{0.34\textwidth}p{0.6\textwidth}}
\hline
Paper item & Lean nodes \\
\hline
Fig.~21, overall garbling scheme & \cref{Shared.programPerfectCorrectness,Shared.perfectCorrectness,RCBComplete.perfectCorrectness,RCBComplete.evaluateCorrect,RCBComplete.evaluateCorrectValid,Garbling.evaluateEncodeRows,Pipeline.evaluateEncoded,RCBComplete.decodeExpectedResult} \\
Eq.~27, $f[r]$ rejects off-curve inputs & \cref{RCBComplete.evaluateCorrectInvalid,RCBComplete.pipelineEvaluateInvalid} \\
Fig.~17 and Fig.~18, circuit $C_4$ & \cref{CurveMembership.evaluateEncodedOnCurve} \\
Claim~7 and Section~5.5, $\mathsf{Enc}_t$ through $\mathsf{EncPRF}$ & \cref{EncPRF.transformEncode} \\
Fig.~9 and Fig.~10, circuit $C_2$ & \cref{FieldMacToECMac.evaluateEncoded,RCBComplete.decodeRowsForOutputKeys,RCBComplete.decodeEvaluateOutputKeyRow} \\
Eqs.~29--31, Jacobian coordinates of $r_j \pi + K_j$ & \cref{RCBComplete.decodeEvaluateRowNone,RCBComplete.decodeEvaluateRowSome,RCBComplete.decodeScaledAlgorithmValue,RCBComplete.outputNonsingular,RCBComplete.outputRepresentsSum} \\
Fig.~6 and Fig.~7, circuit $C_1$ & \cref{Construction.correct,RCBComplete.outputKeyPoints,Construction.offsetsCancel} \\
Appendix~C and Alg.~8, base-$(2-\omega)$ digits & \cref{Construction.scalarReconstruction,GLV91.glvInitialTerminates92} \\
Theorem~8, adaptive privacy of the optimized scheme & \cref{ArithmeticSimulator.lazyCompiledAdaptivePrivacy,Security.sharedStrictOraclePrivacy} \\
Theorem~7, hybrid $\mathsf{Hyb}_0$ & \cref{Security.OperationalOracle.shared_lazy_real} \\
Lemma~14, concrete bound for $\mathsf{GC}_3$ in the pRPM & \cref{Security.sharedStrictWirePrivacy,Security.sharedStrictWireAdvantage_envelope,Security.sharedStrictAdaptiveAdvantage_envelope,Security.sharedStrictGateSource_entrance,Security.sharedAdaptiveThreeRoundingLossSum_le_envelope,Security.sharedErrors_has100Bits,Security.sharedLargeBudget_has100Bits} \\
Lemma~13, H-coefficient technique & \cref{Security.sharedStrictAdaptive_event_bound,Security.hCoefficient_event_of_sourceGoodMass_congr} \\
Claim~11, bad transcript mass $\varepsilon_1$ & \cref{Security.sharedCombinedBad_real_mass_le,Security.sharedFullPipelinePrefixBad_mass_le,Security.sharedGoodMaskSourceBad_mass_le} \\
Claim~12, good transcript ratio & \cref{Security.sharedCombinedGood_real_le,Security.sharedFullPipelinePrefix_real_le_supported,Security.sharedFullCurveGood_real_le_budget} \\
Definition~20, bad transcripts & \cref{Security.sharedStrictTranscript_good,Security.sharedStrictSelectedView_good} \\
Construction~3, simulator $(\mathsf{Sim}_1, \mathsf{Sim}_2)$ & \cref{Security.sharedStrictSourceDecision_allowance,Security.sharedStrictSourceDecision_law,ArithmeticSimulator.lazyCompiledIdealGame_eqStrict,ArithmeticSimulator.lazyCompiledIdealGame_finiteSource,ArithmeticSimulator.lazySetupJoint_parsed,ArithmeticSimulator.lazyChosenDecision_completion,ArithmeticSimulator.lazyCompiledOnline_decision,ArithmeticSimulator.lazyCompiledMachine_onlineParsed,ArithmeticSimulator.lazyOnlineMachine_run,ArithmeticSimulator.lazyCompiledMachine,ArithmeticSimulator.phaseBudget} \\
Eqs.~15--16 and Section~5.1, Lamport labels & \cref{Shared.programLamportCompatible,Lamport.selectedLabels_eq,Lamport.restore_selected,Lamport.keyPairs,Lamport.affineLamportBits_eq_append,Lamport.appendGetLow,Lamport.selectedLabelsEncodeLow,Lamport.keyPairsGetLow,Lamport.appendGetHigh,Lamport.selectedLabelsEncodeHigh,Lamport.keyPairsGetHigh} \\
\hline
\end{tabular}
\caption{Paper items and the Lean nodes that formalize them.}
\label{tab:paper}
\end{table}

\chapter{Perfect correctness}

The root is \cref{Shared.programPerfectCorrectness}.
For every secret scalar $r$, every random tape, and every input $\pi$, $\mathsf{Eval}$ on the garbled circuit and the labels of $\pi$ returns $f[r](\pi)$ (Fig.~21).
The proof follows the evaluation chain of the paper: the curve membership circuit $C_4$ releases $t$, $\mathsf{Dec}_t$ recovers the labels $L_3$, the garbled $C_2$ returns the Jacobian coordinates of $r_j \pi + K_j$, and $\mathsf{Eval}_1$ sums them with the base-$(2-\omega)$ weights.
\Cref{fig:perfect-correctness} shows the proof tree.
An arrow points from a lemma to the result that uses it.

\depgraphfigure{perfect-correctness}{Dependencies of the perfect correctness proof.}

\section{Statements from the paper}

\begin{paperdef}[{Correctness; Definition~3 of the BaBe paper~\cite{babe}}]\label{babe:def3-correctness}
A garbling scheme $\mathsf{GC} = (\mathsf{Garble}, \mathsf{Encode}, \mathsf{Eval})$ is correct if for every circuit $C$ and every input $x$:
\[
  \Pr\big[\, C(x) = \mathsf{Eval}(\mathsf{ct}_{\mathsf{GC}}, L_x) \;:\; (\mathsf{ct}_{\mathsf{GC}}, \mathsf{ek}) \gets \mathsf{Garble}(C),\; L_x \gets \mathsf{Encode}(\mathsf{ek}, x) \,\big] = 1 .
\]
Formalized by \cref{Shared.programPerfectCorrectness}.
\end{paperdef}

\begin{paperdef}[{The function $f[r]$; Eq.~27 of the BaBe paper~\cite{babe}}]\label{babe:eq27}
The Verifier fixes a secret scalar $r \in \mathbb{F}_{\mathbb{r}}^*$.
The Prover supplies a public input point $\pi$ through the bit decomposition $\mathbf{b} = (x_0(\pi), \ldots, x_{n-1}(\pi), y_0(\pi), \ldots, y_{n-1}(\pi))$ of its affine coordinates, with $n = 254$.
The garbled function is
\[
  f[r](\pi) = \begin{cases} r\pi & \text{if } \pi \in \mathbb{G}_1, \\ 0 & \text{otherwise.} \end{cases}
\]
Formalized by \cref{RCBComplete.evaluateCorrect}; the off-curve branch is \cref{RCBComplete.evaluateCorrectInvalid}.
\end{paperdef}

\begin{paperconstr}[{Overall garbling scheme; Fig.~21 and Fig.~29 of the BaBe paper~\cite{babe}}]\label{babe:fig21}
$\mathsf{Garble}(C[r_0, \ldots, r_{\ell-1}])$ runs $\mathsf{Garble}_1(C_1)$, $\mathsf{Garble}_2(C_2[\mathsf{ek}_1])$, $\mathsf{Garble}_3(C_3[\mathsf{ek}_2])$, samples $t, w \gets \mathbb{F}_p$, runs $\mathsf{Garble}_4(C_4[t, w])$ and $\mathsf{Garble}_5(C_5[\mathsf{ek}_4])$, and returns $\mathsf{ct}_{\mathsf{gc}} = (\mathsf{ct}_{\mathsf{gc},2}, \mathsf{ct}_{\mathsf{gc},3}, \mathsf{ct}_{\mathsf{gc},4}, \mathsf{ct}_{\mathsf{gc},5})$ and $\mathsf{ek} = (\mathsf{ek}_3, t, \mathsf{ek}_5)$.
$\mathsf{Encode}(\mathsf{ek}, \mathbf{b})$ returns $L = (\mathsf{Enc}_t(L_3), L_5)$ with $L_3 = \mathsf{Encode}_3(\mathsf{ek}_3, \mathbf{b})$ and $L_5 = \mathsf{Encode}_5(\mathsf{ek}_5, \mathbf{b})$.
$\mathsf{Eval}(\mathsf{ct}_{\mathsf{gc}}, L, \mathbf{b})$ computes $L_4 \gets \mathsf{Eval}_5$, $\hat t \gets \mathsf{Eval}_4(\mathsf{ct}_{\mathsf{gc},4}, L_4, x(\pi), y(\pi))$, $L_3 \gets \mathsf{Dec}_{\hat t}(\mathsf{Enc}_t(L_3))$, $L_2 \gets \mathsf{Eval}_3$, the Jacobian coordinates $(X, Y, Z)(r_j \pi + K_j)_{j} \gets \mathsf{Eval}_2(\mathsf{ct}_{\mathsf{gc},2}, L_2, x(\pi), y(\pi))$, converts them to affine points $L_{1,j}$, and returns $\mathsf{Eval}_1(\bot, L_1)$.
The optimized scheme (Fig.~29) uses the base-$(2-\omega)$ digits $r_j \in D$ with $\ell = 92$ and merges the layers into $\mathsf{GC}_{23}$ and $\mathsf{GC}_{45}$.
Formalized by \cref{RCBComplete.evaluateCorrectValid,Garbling.evaluateEncodeRows,Pipeline.evaluateEncoded,RCBComplete.decodeExpectedResult}.
\end{paperconstr}

\begin{paperconstr}[{Circuit $C_1$ and its garbling; Fig.~6 and Fig.~7 of the BaBe paper~\cite{babe}}]\label{babe:fig7}
$C_1[r_0, \ldots, r_{\ell-1}](\pi) = \sum_{j=0}^{\ell-1} (2-\omega)^j (r_j \pi)$.
$\mathsf{Garble}_1$ samples $K_j \gets \mathbb{G}_1$ for $j \geq 1$, sets $K_0 = -\sum_{j \geq 1} (2-\omega)^j K_j$, and returns $\mathsf{ct}_{\mathsf{gc},1} = \bot$ and $\mathsf{ek}_1 = (r_0, \ldots, r_{\ell-1}, K_0, \ldots, K_{\ell-1})$.
$\mathsf{Encode}_1(\mathsf{ek}_1, \pi)$ returns $L_1 = (L_{1,j})_j$ with $L_{1,j} = r_j \pi + K_j$.
$\mathsf{Eval}_1(\bot, L_1) = \sum_{j=0}^{\ell-1} (2-\omega)^j L_{1,j}$.
Formalized by \cref{Construction.correct,RCBComplete.outputKeyPoints,Construction.offsetsCancel}.
\end{paperconstr}

\begin{paperconstr}[{Circuit $C_2$ and its garbling; Fig.~9, Fig.~10, Eqs.~29--31 of the BaBe paper~\cite{babe}}]\label{babe:fig10}
For each digit $j$ and coordinate $W \in \{X, Y, Z\}$, the Jacobian coordinate $W(r_j \pi + K_j)$ is a quadratic function of $(x(\pi), y(\pi))$:
\[
  W(r_j \pi + K_j) = c^{(j,W)}_0 + c^{(j,W)}_1\, x(\pi) + c^{(j,W)}_2\, y(\pi) + c^{(j,W)}_3\, x(\pi) y(\pi) + c^{(j,W)}_4\, x(\pi)^2 + c^{(j,W)}_5\, y(\pi)^2 .
\]
For a binary digit $r_j$, an offset $K_j = (x(K_j), y(K_j))$, a randomizer $a_j \gets \mathbb{F}_p$, and $B = 3$, the coefficients are
\begin{align*}
  \mathbf{c}^{(j,X)} &= a_j^2 \big( r_j \cdot 2B + (1 - r_j)\, x(K_j),\; r_j\, x(K_j)^2,\; -2 r_j\, y(K_j),\; 0,\; r_j\, x(K_j),\; 0 \big), \\
  \mathbf{c}^{(j,Y)} &= a_j^3 \big( r_j \cdot 3B\, y(K_j) + (1 - r_j)\, y(K_j),\; 0,\; r_j (-y(K_j)^2 - 3B),\; -3 r_j\, x(K_j)^2,\; 3 r_j\, x(K_j) y(K_j),\; r_j\, y(K_j) \big), \\
  \mathbf{c}^{(j,Z)} &= a_j \big( -r_j\, x(K_j) + (1 - r_j),\; r_j,\; 0,\; 0,\; 0,\; 0 \big).
\end{align*}
The base-$(2-\omega)$ implementation uses the adjusted coefficients of \cite[Eqs.~51--53]{babe} for digits $r_j \in D$.
$\mathsf{Garble}_2$ applies the Ishai--Wee partial garbling: it samples masks $\rho^{(j,W)}_1, \ldots, \rho^{(j,W)}_{10}$, sets $\rho_0 = -\rho_9 - \rho_{10}$, and publishes $c'_k = c_k + \rho_k$ for $k \leq 5$.
$\mathsf{Encode}_2$ returns the labels $c'_6, \ldots, c'_{10}$, which are affine in $x(\pi)$ and $y(\pi)$.
$\mathsf{Eval}_2$ returns $c'_0 + c'_1 x + c'_2 y + c'_3 xy + c'_4 x^2 + c'_5 y^2 + c'_6 x + c'_7 x + c'_8 y + c'_9 + c'_{10} = W(r_j \pi + K_j)$.
Formalized by \cref{FieldMacToECMac.evaluateEncoded,RCBComplete.decodeRowsForOutputKeys,RCBComplete.decodeEvaluateOutputKeyRow,RCBComplete.decodeEvaluateRowNone,RCBComplete.decodeEvaluateRowSome,RCBComplete.decodeScaledAlgorithmValue,RCBComplete.outputNonsingular,RCBComplete.outputRepresentsSum}.
\end{paperconstr}

\begin{paperconstr}[{Curve membership circuit $C_4$; Fig.~17 and Fig.~18 of the BaBe paper~\cite{babe}}]\label{babe:fig18}
$C_4[t, w](x, y) = t + w\,(x^3 + 3 - y^2) = (t + 3w) + w\, x^3 - w\, y^2$ for $t, w \gets \mathbb{F}_p$.
The value $t$ is released if and only if $\pi$ is on the curve.
$\mathsf{Garble}_4$ applies the Ishai--Wee technique to the coefficients $(\hat c_0, \hat c_1, \hat c_2) = (t + 3w, w, -w)$ with masks $\hat\rho_1, \ldots, \hat\rho_7$, and $\mathsf{Eval}_4$ returns $\hat c'_0 + \hat c'_1 x^3 + \hat c'_2 y^2 + \hat c'_3 x^2 + \hat c'_4 y + \hat c'_5 x + \hat c'_6 + \hat c'_7$.
Formalized by \cref{CurveMembership.evaluateEncodedOnCurve}.
\end{paperconstr}

\begin{paperclaim}[{Encryption of the labels; Claim~7 and Eq.~33 of the BaBe paper~\cite{babe}}]\label{babe:claim7}
$(\mathsf{Enc}, \mathsf{Dec})$ is implemented with a PRF $\mathsf{EncPRF}$ and a public counter: $\mathsf{Enc}_t(m) = \mathsf{EncPRF}(t, \mathsf{counter}) \oplus m$ and $\mathsf{Dec}_t(\mathsf{ct}) = \mathsf{EncPRF}(t, \mathsf{counter}) \oplus \mathsf{ct}$.
$\mathsf{EncPRF}(t, (w, i, b)) = \mathsf{P}_{w,i}(\mathsf{enc}(b) \oplus k_1) \oplus k_2$ with $k_1 \| k_2 = \mathsf{SHA256}(t)$ is the two-key Even--Mansour construction, where each $\mathsf{P}_{w,i}$ is a programmable random permutation.
The labels of $\mathsf{GC}_3$ are encrypted label by label: $\hat L^b_{x,i} = \mathsf{Enc}_t(L^b_{x,i})$ and $\hat L^b_{y,i} = \mathsf{Enc}_t(L^b_{y,i})$ (Eq.~33).
If $\mathsf{GC}_{45}$ is $\epsilon_1$-adaptively private, then the first two hybrids of the privacy proof differ by at most $\epsilon_1 + \mathsf{negl}(\lambda)$.
Formalized by \cref{EncPRF.transformEncode}.
\end{paperclaim}

\begin{paperalg}[{Base-$(2-\omega)$ digits; Appendix~C and Alg.~8 of the BaBe paper~\cite{babe}}]\label{babe:alg8}
Let $\omega \in \mathbb{F}_{\mathbb{r}}$ with $\omega^2 + \omega + 1 = 0$ and $D = \{0, \pm 1, \pm\omega, \pm(1 + \omega)\}$.
For every scalar $r \in \{0, \ldots, \mathbb{r} - 1\}$, the GLV decomposition gives small integers $(a, b)$ with $r \equiv a + \omega b \pmod{\mathbb{r}}$.
Alg.~8 divides $A + B\omega$ by $2 - \omega$ repeatedly, with the digit $r_i \in D$ chosen from $(A + 2B) \bmod 7$ and the exact division $A' = (3A - B)/7$, $B' = (A + 2B)/7$, and stops at $(A, B) = (0, 0)$ within $\ell = 1 + \lceil \log_7 \mathbb{r} \rceil = 92$ iterations.
The reconstruction $\sum_{i=0}^{\ell-1} r_i (2-\omega)^i \equiv r \pmod{\mathbb{r}}$ holds.
Multiplication by a digit in $D$ is an endomorphism: $\omega \pi = (\mu\, x(\pi), y(\pi))$ with $\mu^2 + \mu + 1 = 0$ in $\mathbb{F}_p$, and $-\pi = (x(\pi), -y(\pi))$.
Formalized by \cref{Construction.scalarReconstruction,GLV91.glvInitialTerminates92}.
\end{paperalg}

\section{Lean statements}

\subsection{Program and wire circuits}
\inputleannode{Shared.programPerfectCorrectness}
\inputleannode{Shared.perfectCorrectness}
\inputleannode{Lamport.restore_selected}
\inputleannode{RCBComplete.perfectCorrectness}
\inputleannode{RCBComplete.evaluateCorrect}

\subsection{Invalid inputs}
\inputleannode{RCBComplete.evaluateCorrectInvalid}
\inputleannode{RCBComplete.pipelineEvaluateInvalid}

\subsection{Valid inputs}
\inputleannode{RCBComplete.evaluateCorrectValid}
\inputleannode{Garbling.evaluateEncodeRows}
\inputleannode{Pipeline.evaluateEncoded}
\inputleannode{CurveMembership.evaluateEncodedOnCurve}
\inputleannode{EncPRF.transformEncode}
\inputleannode{FieldMacToECMac.evaluateEncoded}

\subsection{Decoding}
\inputleannode{RCBComplete.decodeExpectedResult}
\inputleannode{RCBComplete.decodeRowsForOutputKeys}
\inputleannode{RCBComplete.decodeEvaluateOutputKeyRow}
\inputleannode{RCBComplete.decodeEvaluateRowNone}
\inputleannode{RCBComplete.decodeEvaluateRowSome}
\inputleannode{RCBComplete.decodeScaledAlgorithmValue}
\inputleannode{RCBComplete.outputNonsingular}
\inputleannode{RCBComplete.outputRepresentsSum}
\inputleannode{RCBComplete.outputKeyPoints}

\subsection{Randomized encoding}
\inputleannode{Construction.correct}
\inputleannode{Construction.scalarReconstruction}
\inputleannode{GLV91.glvInitialTerminates92}
\inputleannode{Construction.offsetsCancel}

\chapter{Adaptive privacy}

The root is \cref{ArithmeticSimulator.lazyCompiledAdaptivePrivacy}.
The requirement is the work-per-advantage level of Lemma~14: every adversary with work $T$ has advantage $\delta$ with $T / \delta \geq 2^{100}$.
The proof follows the H-coefficient technique of Lemma~13 in the programmable random permutation model, with the bad transcripts of Definition~20, the bad-transcript mass $\varepsilon_1$ of Claim~11, and the good-transcript ratio of Claim~12.
The simulator $(\mathsf{Sim}_1, \mathsf{Sim}_2)$ of Construction~3 is compiled to a bounded machine.
\Cref{fig:adaptive-privacy} shows the proof tree.
Grey nodes belong to the perfect correctness chapter.

\depgraphfigurewide{adaptive-privacy}{Dependencies of the adaptive privacy proof.}

\section{Statements from the paper}

\begin{paperdef}[{Adaptive privacy; Definition~3 of the BaBe paper~\cite{babe}}]\label{babe:def3-privacy}
A garbling scheme is $\epsilon$-adaptively private if there is a simulator $\mathsf{Sim} = (\mathsf{Sim}_1, \mathsf{Sim}_2)$ such that for every adversary $\mathcal{A}$, every circuit $C$, and every auxiliary input $\mathsf{aux}$:
\begin{align*}
  \Big| &\Pr\big[ \mathcal{A}(\mathsf{ct}_{\mathsf{GC}}, L_x, \mathsf{aux}) = 1 : (\mathsf{ct}_{\mathsf{GC}}, \mathsf{ek}) \gets \mathsf{Garble}(C),\; x \gets \mathcal{A}(\mathsf{ct}_{\mathsf{GC}}, \mathsf{aux}),\; L_x \gets \mathsf{Encode}(\mathsf{ek}, x) \big] \\
  &- \Pr\big[ \mathcal{A}(\mathsf{ct}_{\mathsf{GC}}, L_x, \mathsf{aux}) = 1 : (\mathsf{ct}_{\mathsf{GC}}, \mathsf{st}) \gets \mathsf{Sim}_1(\mathsf{topo}(C)),\; x \gets \mathcal{A}(\mathsf{ct}_{\mathsf{GC}}, \mathsf{aux}),\; L_x \gets \mathsf{Sim}_2(\mathsf{st}, C(x)) \big] \Big| \leq \epsilon .
\end{align*}
The adversary chooses the input $x$ after it sees the garbled circuit.
The Lean statement bounds the advantage by the work of the adversary: $T / \delta \geq 2^{100}$ with $T = q + 1$.
Formalized by \cref{ArithmeticSimulator.lazyCompiledAdaptivePrivacy,Security.sharedStrictOraclePrivacy}.
\end{paperdef}

\begin{paperthm}[{Theorem~7 of the BaBe paper~\cite{babe}}]\label{babe:thm7}
Suppose $\mathsf{GC}_3$ is $\epsilon_3$-adaptively private and $\mathsf{GC}_5$ is $\epsilon_5$-adaptively private.
Suppose the PRF $\mathsf{CTPRF}$ used in $\mathsf{GC}_3$ and $\mathsf{GC}_5$ is modeled as a programmable random permutation and $(\mathsf{Enc}, \mathsf{Dec})$ is $\epsilon_2$-CPA secure.
Then the garbling scheme $\mathsf{GC}$ of Fig.~21 is $(\epsilon_2 + \epsilon_3 + \epsilon_5)$-adaptively private.
The proof is a sequence of hybrids; $\mathsf{Hyb}_0$ is the real experiment.
Formalized by \cref{Security.OperationalOracle.shared_lazy_real} for $\mathsf{Hyb}_0$.
\end{paperthm}

\begin{paperthm}[{Theorem~8 of the BaBe paper~\cite{babe}}]\label{babe:thm8}
Suppose the PRF $\mathsf{CTPRF}$ used in $\mathsf{GC}_{23}$ and $\mathsf{GC}_{45}$ is modeled as a programmable random permutation.
Suppose $(\mathsf{Enc}, \mathsf{Dec})$ is $\epsilon_2$-CPA secure and is implemented with $\mathsf{EncPRF}$ as in \cref{babe:claim7}, where each permutation $\mathsf{P}_{w,i}$ is a programmable random permutation and $\mathsf{SHA256}$ is a random oracle.
Then the optimized garbling scheme $\mathsf{GC}^{\mathrm{opt}}$ of Fig.~29 satisfies $(\epsilon_2 + \mathsf{negl}(\lambda))$-adaptive privacy.
Formalized by \cref{ArithmeticSimulator.lazyCompiledAdaptivePrivacy,Security.sharedStrictOraclePrivacy}.
\end{paperthm}

\begin{paperlem}[{Key insight; Lemma~12 of the BaBe paper~\cite{babe}}]\label{babe:lem12}
Fix a deterministic adversary.
For every attainable transcript $\tau$ and every world $b \in \{\mathrm{real}, \mathrm{ideal}\}$: $\Pr_{\omega \gets \Omega_b}[Z_b(\omega) = \tau] = \Pr_{\omega \gets \Omega_b}[\omega \in \mathsf{comp}_b(\tau)]$, where $Z_b$ is the transcript of the interaction with the oracle $\omega$ and $\mathsf{comp}_b(\tau)$ is the set of oracles compatible with $\tau$.
\end{paperlem}

\begin{paperlem}[{H-coefficient technique; Lemma~13 of the BaBe paper~\cite{babe}}]\label{babe:lem13}
Let $\mathcal{T}_{\mathrm{bad}}$ be a set of bad transcripts and $\mathcal{T}_{\mathrm{good}}$ its complement in the attainable transcripts.
Suppose there are $\varepsilon_1, \varepsilon_2 \geq 0$ such that
$\sum_{\tau \in \mathcal{T}_{\mathrm{bad}}} \Pr_{\omega \gets \Omega_{\mathrm{ideal}}}[\omega \in \mathsf{comp}_{\mathrm{ideal}}(\tau)] \leq \varepsilon_1$ and, for every $\tau \in \mathcal{T}_{\mathrm{good}}$,
$1 - \Pr_{\omega \gets \Omega_{\mathrm{real}}}[\omega \in \mathsf{comp}_{\mathrm{real}}(\tau)] / \Pr_{\omega \gets \Omega_{\mathrm{ideal}}}[\omega \in \mathsf{comp}_{\mathrm{ideal}}(\tau)] \leq \varepsilon_2$.
Then the advantage of the adversary is at most $\varepsilon_1 + \varepsilon_2$.
Formalized by \cref{Security.hCoefficient_event_of_sourceGoodMass_congr,Security.sharedStrictAdaptive_event_bound}.
\end{paperlem}

\begin{paperlem}[{Adaptive privacy of $\mathsf{GC}_3$ in the pRPM; Lemma~14 of the BaBe paper~\cite{babe} and the remark after it}]\label{babe:lem14}
Let $\mathsf{CTPRF}$ be instantiated from fixed-key AES-128 with one independent permutation per bucket $(i, k, W, \mathsf{slot})$, each modeled as a programmable random permutation at $\lambda = 128$.
Let $n = 254$, let $B = 6858$ be the number of permutations, let $Q^{(b)}_{\mathrm{act}}$ be the number of active gates per bucket, and let $q$ be the total number of permutation queries of the adversary.
Then $\mathsf{GC}_3$ (Construction~2) with the simulator of Construction~3 is $\varepsilon_3$-adaptively private in the pRPM with
\[
  \varepsilon_3 \leq \frac{3 B (Q^{(b)}_{\mathrm{act}})^2 + 4 Q^{(b)}_{\mathrm{act}}\, q}{2^{\lambda}} .
\]
With the binary digits, $Q^{(b)}_{\mathrm{act}} = n = 254$ and the work-per-advantage level $\min_{\mathcal{A}} \log_2(T_{\mathcal{A}} / \delta_{\mathcal{A}})$ is at least $97.6$ bits for every $q$.
With the base-$(2-\omega)$ digits of the implementation, $Q^{(b)}_{\mathrm{act}} = \ell = 92$ and the level is at least $100$ bits: $T / \delta \geq 2^{100.62}$.
Formalized by \cref{Security.sharedStrictWirePrivacy,Security.sharedStrictWireAdvantage_envelope,Security.sharedStrictAdaptiveAdvantage_envelope,Security.sharedAdaptiveThreeRoundingLossSum_le_envelope,Security.sharedErrors_has100Bits,Security.sharedLargeBudget_has100Bits}, with the Lean envelope $\varepsilon_3(q) = (251850146 + 833\,q) / 2^{128}$.
\end{paperlem}

\begin{paperdef}[{Bad transcripts; Definition~20 of the BaBe paper~\cite{babe}}]\label{babe:def20}
For a transcript $\tau = (K_1, \mathsf{ct}_{\mathsf{GC},3}, K_2, (\mathbf{b}, L_3), K_3)$, let $a^*_{i,j,k,W} = L^{(i)}_{3,\delta(k)} \oplus \mathsf{enc}(j, k, W)$ be the active permutation input of gate $(i, j, k, W)$, let $v^*_{i,j,k,W}$ be its required output, and let $w_{i,j,k,W}$ be its inactive output offset.
A transcript is bad if one of the following events occurs.
\begin{description}
  \item[bad1] Two distinct gates have the same active input $a^*$.
  \item[bad2] Two distinct gates have the same required output $v^*$.
  \item[bad3] A permutation query $(\alpha, \beta) \in K_1 \cup K_2$ made before the labels are released hits an active input, $\alpha = a^*$, or a required output, $\beta = v^*$.
  \item[bad4] Two distinct gates that share a wire label have the same inactive offset $w$.
\end{description}
Formalized by \cref{Security.sharedStrictTranscript_good,Security.sharedStrictSelectedView_good}: the strict simulator aborts exactly on a bad sample.
\end{paperdef}

\begin{paperclaim}[{Bad-transcript mass; Claim~11 of the BaBe paper~\cite{babe}}]\label{babe:claim11}
$\varepsilon_1 = \sum_{\tau \in \mathcal{T}_{\mathrm{bad}}} \Pr_{\omega \gets \Omega_{\mathrm{ideal}}}[\omega \in \mathsf{comp}_{\mathrm{ideal}}(\tau)] \leq \dfrac{\tfrac{3}{2} Q_{\mathrm{act}}^2 + 2 q\, Q_{\mathrm{act}}}{2^{\lambda}}$,
where bad1, bad2, and bad4 each contribute at most $Q_{\mathrm{act}}^2 / 2^{\lambda + 1}$ and bad3 contributes at most $2 q\, Q_{\mathrm{act}} / 2^{\lambda}$.
Formalized by \cref{Security.sharedCombinedBad_real_mass_le,Security.sharedFullPipelinePrefixBad_mass_le,Security.sharedGoodMaskSourceBad_mass_le}.
\end{paperclaim}

\begin{paperclaim}[{Good-transcript ratio; Claim~12 of the BaBe paper~\cite{babe}}]\label{babe:claim12}
For every good transcript $\tau$:
$\Pr_{\omega \gets \Omega_{\mathrm{real}}}[\omega \in \mathsf{comp}_{\mathrm{real}}(\tau)] / \Pr_{\omega \gets \Omega_{\mathrm{ideal}}}[\omega \in \mathsf{comp}_{\mathrm{ideal}}(\tau)] \geq 1 - \delta_2$ with $\delta_2 = (3 Q_{\mathrm{act}}^2 + 2 q\, Q_{\mathrm{act}}) / 2^{\lambda}$, so $\varepsilon_2 = \delta_2$.
Formalized by \cref{Security.sharedCombinedGood_real_le,Security.sharedFullPipelinePrefix_real_le_supported,Security.sharedFullCurveGood_real_le_budget}.
\end{paperclaim}

\begin{paperconstr}[{Simulator for $\mathsf{GC}_3$; Construction~3 of the BaBe paper~\cite{babe}}]\label{babe:con3}
$\mathsf{Sim}_1(\mathsf{topo}(C_3))$ samples every ciphertext $\mathsf{ct}^{(i,j,k,W)}_{3,b} \gets \{0,1\}^{\lambda}$ uniformly and outputs $\mathsf{ct}_{\mathsf{GC},3}$ as its state.
$\mathsf{Sim}_2(\mathsf{st}, \mathbf{b}, \{c'^{(j,W)}_{k,i}\})$ samples the active labels $L^{(i)}_{3,x}, L^{(i)}_{3,y} \gets \{0,1\}^{\lambda}$, and for every gate programs the permutation at $L^{(i)}_{3,\delta(k)} \oplus \mathsf{enc}(j, k, W)$ so that $\mathsf{CTPRF}$ decrypts the active ciphertext to the correct output $c'^{(j,W)}_{k,i}$; it outputs $L_3$.
The Lean simulator compiles $(\mathsf{Sim}_1, \mathsf{Sim}_2)$ to a bounded machine with at most $2^{60}$ instructions, and bounds every rejection-sampling draw to $256$ attempts.
Formalized by \cref{ArithmeticSimulator.lazyCompiledMachine,ArithmeticSimulator.lazyCompiledIdealGame_eqStrict,Security.sharedStrictSourceDecision_law,Security.sharedStrictSourceDecision_allowance}.
\end{paperconstr}

\section{Lean statements}

\subsection{Closed simulator}
\inputleannode{ArithmeticSimulator.lazyCompiledAdaptivePrivacy}
\inputleannode{Security.sharedStrictOraclePrivacy}
\inputleannode{Security.OperationalOracle.shared_lazy_real}
\inputleannode{Security.sharedStrictWirePrivacy}

\subsection{Privacy bound}
\inputleannode{Security.sharedStrictWireAdvantage_envelope}
\inputleannode{Security.sharedStrictAdaptiveAdvantage_envelope}
\inputleannode{Security.sharedStrictAdaptive_event_bound}
\inputleannode{Security.hCoefficient_event_of_sourceGoodMass_congr}
\inputleannode{Security.sharedCombinedBad_real_mass_le}
\inputleannode{Security.sharedCombinedGood_real_le}
\inputleannode{Security.sharedFullPipelinePrefix_real_le_supported}
\inputleannode{Security.sharedFullCurveGood_real_le_budget}
\inputleannode{Security.sharedStrictTranscript_good}
\inputleannode{Security.sharedStrictSelectedView_good}
\inputleannode{Security.sharedStrictGateSource_entrance}

\subsection{Collision bounds}
\inputleannode{Security.sharedFullPipelinePrefixBad_mass_le}
\inputleannode{Security.sharedGoodMaskSourceBad_mass_le}
\inputleannode{Security.sharedAdaptiveThreeRoundingLossSum_le_envelope}

\subsection{Finite sampling}
\inputleannode{Security.sharedStrictSourceDecision_allowance}
\inputleannode{Security.sharedStrictSourceDecision_law}
\inputleannode{Security.sharedErrors_has100Bits}
\inputleannode{Security.sharedLargeBudget_has100Bits}

\subsection{Simulator execution}
\inputleannode{ArithmeticSimulator.lazyCompiledIdealGame_eqStrict}
\inputleannode{ArithmeticSimulator.lazyCompiledIdealGame_finiteSource}
\inputleannode{ArithmeticSimulator.lazySetupJoint_parsed}
\inputleannode{ArithmeticSimulator.lazyChosenDecision_completion}
\inputleannode{ArithmeticSimulator.lazyCompiledOnline_decision}
\inputleannode{ArithmeticSimulator.lazyCompiledMachine_onlineParsed}
\inputleannode{ArithmeticSimulator.lazyOnlineMachine_run}
\inputleannode{ArithmeticSimulator.lazyCompiledMachine}
\inputleannode{ArithmeticSimulator.phaseBudget}

\chapter{Lamport compatibility}

The root is \cref{Shared.programLamportCompatible}.
The input labels $\mathsf{Encode}(\mathsf{ek}, \pi)$ are exactly the Lamport signature of the bits of $\pi$ under the key pairs of $\mathsf{ek}$, as the locking scripts of Eqs.~15--16 require.
\Cref{fig:lamport-compatibility} shows the proof tree.

\depgraphfigure{lamport-compatibility}{Dependencies of the Lamport compatibility proof.}

\section{Statements from the paper}

\begin{paperdef}[{Bitcoin constraint on the input labels; Section~5.1 of the BaBe paper~\cite{babe}}]\label{babe:bitcoin}
The input labels produced by $\mathsf{Encode}$ are revealed on-chain.
The Prover commits to $\pi$ by a Lamport signature on the bits $\mathbf{b}$ of its affine coordinates, and the input labels of the garbled circuit are exactly the Lamport signatures of $\mathbf{b}$ under the encoding key $\mathsf{ek}$: for the bit $x_i(\pi)$ the label is $L^{x_i(\pi)}_{x,i}$ and for the bit $y_i(\pi)$ the label is $L^{y_i(\pi)}_{y,i}$.
Formalized by \cref{Shared.programLamportCompatible,Lamport.selectedLabels_eq}.
\end{paperdef}

\begin{paperdef}[{Locking scripts; Eqs.~15--16 of the BaBe paper~\cite{babe}}]\label{babe:eq15}
For a Lamport public key $\mathsf{lpk} = (\mathsf{lpk}^0_i, \mathsf{lpk}^1_i)_{i \leq \ell}$, the script
$\mathsf{CheckLampSig}(\mathsf{lpk}) = \bigwedge_{i=1}^{\ell} \big( \mathsf{HashLock}(\mathsf{lpk}^0_i) \lor \mathsf{HashLock}(\mathsf{lpk}^1_i) \big)$
requires a Lamport signature on some message, and
$\mathsf{CheckLampSigsMatch}(\mathsf{lpk}_A, \mathsf{lpk}_B) = \bigwedge_{i=1}^{\ell} \big[ (\mathsf{HashLock}((\mathsf{lpk}_A)^0_i) \land \mathsf{HashLock}((\mathsf{lpk}_B)^0_i)) \lor (\mathsf{HashLock}((\mathsf{lpk}_A)^1_i) \land \mathsf{HashLock}((\mathsf{lpk}_B)^1_i)) \big]$
requires signatures under both keys on the same message.
In ArgoMAC the message is $\mathbf{b}$ with $\ell = 508$ bits, and the key pairs are the labels $(L^0_{x,i}, L^1_{x,i})$ and $(L^0_{y,i}, L^1_{y,i})$.
Formalized by \cref{Lamport.keyPairs,Lamport.restore_selected}.
\end{paperdef}

\section{Lean statements}

\inputleannode{Shared.programLamportCompatible}
\inputleannode{Lamport.keyPairs}
\inputleannode{Lamport.selectedLabels_eq}
\inputleannode{Lamport.affineLamportBits_eq_append}
\inputleannode{Lamport.appendGetLow}
\inputleannode{Lamport.selectedLabelsEncodeLow}
\inputleannode{Lamport.keyPairsGetLow}
\inputleannode{Lamport.appendGetHigh}
\inputleannode{Lamport.selectedLabelsEncodeHigh}
\inputleannode{Lamport.keyPairsGetHigh}
